\subsection{Methodology}
\label{sec:method-P}

In this section we reproduce Section~3.2 of \citet{darcet2024vision}, which studies whether adding register tokens to a frozen ViT degrades or improves downstream linear evaluation. We compare three backbones: \textbf{DINOv2 ViT-L/14} (24 transformer blocks, embedding dimension $1024$, 16 attention heads, $\sim 304$M parameters), \textbf{OpenCLIP ViT-B/16} (12 blocks, dimension $768$, 12 heads, $\sim 86$M parameters in the vision tower), and \textbf{DeiT-III ViT-B/16} (12 blocks, dimension $768$, 12 heads, $\sim 87$M parameters). Each is tested in a no-register and a with-register variant. Since no trained-register checkpoint is publicly available for OpenCLIP or DeiT-III, the with-register variants are substituted by the test-time register method of \citet{jiang2025vision} for OpenCLIP and by injecting 4 untrained register-token \texttt{nn.Parameter}s at inference for DeiT-III (both are documented limitations). For the Figure~8 ablation we additionally compare two further DINOv2-based interventions (truncate/cycle of the released reg4 checkpoint, and an independent variant of the Jiang method). All code, including the from-scratch re-implementation of \citet{jiang2025vision} Algorithm~1 for the \texttt{timm} GELU port of DINOv2, is in \texttt{Adv\_ML/src/Paolo/}.

\subsubsection{Datasets}
\citet{darcet2024vision} report Table~2a on the full versions of three benchmarks: \textbf{ImageNet-1k} (1000 classes, $\sim$1.28M training images, 50K validation images), \textbf{ADE20k} (150 semantic classes plus a background label, $\sim$20K training images, $\sim$2K validation images), and \textbf{NYU Depth v2} (indoor scenes with Kinect-acquired RGB+depth pairs, $\sim$24K training images, $\sim$700 validation images). Training the linear heads on the full training sets is not feasible in our compute budget; we therefore use stratified subsets of 500 images each, sampled from the \emph{official validation split} of every dataset and split 7:3 into 350 train / 150 val with no separate test set. Concretely: for ImageNet we sample 50 classes out of 1000 (seed~42) and take 10 images per class; for ADE20k and NYUd we sample 500 images uniformly at random because their per-image content is multi-class (segmentation) or continuous (depth), so stratification by class is not meaningful. We do not need a separate test set because the only hyperparameter the linear heads expose (sklearn's \texttt{C}) is kept at its default and never tuned, so there is no risk of validation-set leakage. As a consequence of this much smaller pool, our absolute numbers do not match the paper's full-dataset numbers; we compare only the \emph{relative} effect of adding registers. Dataset references are in the shared Table~\ref{tab:datasets}.

\subsubsection{Hyperparameters}
The classification head is \texttt{sklearn.LogisticRegression(max\_iter=1000)} fitted on frozen CLS features with all other defaults. The segmentation head is a $1\!\times\!1$ \texttt{Conv2d}($C \to 151$) with cross-entropy loss (ignore-index~0), AdamW ($\text{lr}=10^{-3}$, weight-decay~$10^{-2}$), batch size~4, 30~epochs. The depth head is identical except $\text{Conv2d}(C \to 1)$ trained on log-depth with MSE and the output clamped to the dataset depth range before exponentiation. For the test-time register intervention we use the hyperparameters specified in Jiang's \texttt{configs/dinov2\_large.yaml}: \texttt{top\_k}~$=50$ register neurons, \texttt{register\_norm\_threshold}~$=150$, neuron identification on 100 ImageNet training images.

\subsubsection{Experimental setup and code}
We define the metrics for reproducibility: \textbf{Top-1 accuracy} is the fraction of validation images whose predicted class equals the ground truth; \textbf{mIoU} is the mean over 150 ADE20k classes of $\mathrm{IoU} = \mathrm{TP}/(\mathrm{TP}+\mathrm{FP}+\mathrm{FN})$ at the pixel level; \textbf{RMSE} is $\sqrt{\tfrac{1}{N}\sum (\hat{d}_i - d_i)^2}$ in metres on NYUd. For Table~2a we run the three probes on each $\{$backbone$\}\times\{$no-reg, with-reg$\}$ combination; the model checkpoints used are listed in the shared models table (Table~\ref{tab:models}). For Table~2b we run zero-shot ImageNet on OpenCLIP using the standard 11-prompt template ensemble averaged in CLIP's joint embedding space. For Figure~8 we sweep $N \in \{0,1,2,4,8,16\}$ register tokens using three approximations to a full retraining sweep (which is out of scope for our compute budget): \textbf{(1)~truncate/cycle} on DINOv2-reg4 -- for $N<4$ we truncate the four trained registers to the first $N$ slots; for $N>4$ we cycle them to fill the additional slots; \textbf{(2)~test-time registers} \citep{jiang2025vision} on the no-register DINOv2 baseline -- we re-implement Algorithm~1 (\textsc{FindRegisterNeurons}) and the forward-hook intervention from scratch on the \texttt{timm} GELU port, and validate by confirming that five of our top-10 register neurons coincide with the indices the authors publish for their Meta-DINOv2 SwiGLU port; \textbf{(3)~independent variant} (ours) -- identical pipeline to (2), but for each neuron the top-$N$ distinct outlier values are written onto the $N$ test-time registers rather than the top-1 value duplicated $N$ times. Attention maps for Figure~8 (top) are extracted from the last self-attention layer of DINOv2 with \texttt{fused\_attn} disabled so that explicit softmax probabilities can be captured by a forward hook.

\subsubsection{Computational requirements}
All experiments ran on an Apple M3 chip via the PyTorch MPS backend; no dedicated GPU was used. Rounded wall-clock estimates per elementary operation (feature extraction, head training, register-neuron identification, single-image inference) and per top-level experiment (Tables~2a/2b, three Figure~8 sweeps) are reported in Table~\ref{tab:timings}. The total wall-clock budget for the EXP2 experiments is estimated at $\sim$12--14~hours; the dominant cost is repeated feature extraction with DINOv2 ViT-L/14 across the three Figure~8 sweeps.
